\section{학습 및 평가}

\subsection{학습 단계}

기반 모델은 Qwen3.5-9B이다. 모든 PT에서 전체 파라미터를 1 epoch
학습하며 bfloat16, cosine learning rate $10^{-5}$, warmup 3\%,
completion-only loss, 최대 길이 8192, effective batch size 256을
사용한다. 이미지는 최대 2,359,296 pixel로 입력하며 현재 해상도에서
3,009 image token에 해당한다. 모든 PT run에 gradient checkpointing을
적용한다.

PT 이후 Base, Short PT, Long PT 모델은 같은 순서의 100K IT 예제를
학습한다. 이 데이터는 MolmoWeb SyntheticQA 70K와 필터링한
LLaVA-NeXT 30K로 구성한다. 후속 일반화 평가와의 중복을 막기 위해
LLaVA-NeXT에서 DocVQA, ChartQA, DVQA, AI2D 출처를 제거한다. 세 IT
체크포인트에는 다시 동일한 99,999개의 trajectory 예제를 학습한다.
행동 공간은 \texttt{click}, \texttt{type}, \texttt{press},
\texttt{scroll}, \texttt{goto}, \texttt{go\_back}, \texttt{finish}의
일곱 가지로 통일한다.

\begin{figure*}[t]
\centering
\begin{tikzpicture}[
  font=\small,
  stage/.style={draw, rounded corners=2pt, align=center,
                minimum width=20mm, minimum height=8mm},
  skip/.style={draw, rounded corners=2pt, align=center,
               minimum width=20mm, minimum height=8mm, dashed,
               text=gray},
  arrow/.style={-{Latex[length=2mm]}, semithick},
  row/.style={font=\small\bfseries, anchor=east}
]
\node[row] at (-0.5,0) {No PT};
\node[row] at (-0.5,-1.05) {PT, no IT};
\node[row] at (-0.5,-2.10) {Short};
\node[row] at (-0.5,-3.15) {Long bilingual};

\node[stage,fill=stagegray] (base0) at (1,0) {Base};
\node[stage,fill=stagegray] (base1) at (1,-1.05) {Base};
\node[stage,fill=stagegray] (base2) at (1,-2.10) {Base};
\node[stage,fill=stagegray] (base3) at (1,-3.15) {Base};

\node[skip]  (pt0) at (3.6,0) {PT};
\node[stage,fill=stagegray] (it0) at (6.2,0) {IT};
\node[stage,fill=stagegreen] (act0) at (8.8,0) {Action};
\node[stage,fill=stagegreen] (traj0) at (11.4,0) {Trajectory};

\node[stage,fill=longorange] (pt1) at (3.6,-1.05) {PT};
\node[skip]  (it1) at (6.2,-1.05) {IT};
\node[stage,fill=stagegreen] (act1) at (8.8,-1.05) {Action};
\node[stage,fill=stagegreen] (traj1) at (11.4,-1.05) {Trajectory};

\node[stage,fill=shortblue] (pt2) at (3.6,-2.10) {Short PT};
\node[stage,fill=stagegray] (it2) at (6.2,-2.10) {IT};
\node[stage,fill=stagegreen] (act2) at (8.8,-2.10) {Action};
\node[stage,fill=stagegreen] (traj2) at (11.4,-2.10) {Trajectory};

\node[stage,fill=longorange] (pt3) at (3.6,-3.15) {Long PT\\EN+KO};
\node[stage,fill=stagegray] (it3) at (6.2,-3.15) {IT};
\node[stage,fill=stagegreen] (act3) at (8.8,-3.15) {Action};
\node[stage,fill=stagegreen] (traj3) at (11.4,-3.15) {Trajectory};

\foreach \r in {0,1,2,3} {
  \draw[arrow] (base\r) -- (pt\r);
  \draw[arrow] (pt\r) -- (it\r);
  \draw[arrow] (it\r) -- (act\r);
  \draw[arrow] (act\r) -- (traj\r);
}
\end{tikzpicture}
\caption{\textbf{실제로 준비한 네 학습 경로.} 점선 상자는 해당
단계를 건너뜀을 뜻한다. 각 단계의 데이터 구성과 비교 목적은
본문에서 설명한다.}
\label{fig:pipeline}
\end{figure*}


그림 \ref{fig:pipeline}과 표 \ref{tab:checkpoints}는 전체 비교 구조를
보여준다. PT가 없는 분기는 PT 자체의 효과와 Short/Long 간 차이를
분리하는 기준선이다.

\begin{table}[t]
\centering
\small
\setlength{\tabcolsep}{3.5pt}
\begin{tabular}{clccc}
\toprule
ID & 체크포인트 & PT & IT & Traj \\
\midrule
C0 & Base & -- & -- & -- \\
C1 & Short PT & Short & -- & -- \\
C2 & Long PT & Long & -- & -- \\
C3 & Base + IT & -- & \checkmark & -- \\
C4 & Short PT + IT & Short & \checkmark & -- \\
C5 & Long PT + IT & Long & \checkmark & -- \\
C6 & Base + IT + Traj & -- & \checkmark & \checkmark \\
C7 & Short + IT + Traj & Short & \checkmark & \checkmark \\
C8 & Long + IT + Traj & Long & \checkmark & \checkmark \\
\bottomrule
\end{tabular}
\caption{\textbf{평가할 아홉 개 체크포인트.} 중간 체크포인트를
보존하여 PT 효과가 나타나는 지점과 후속 학습 이후의 유지 여부를
분리한다.}
\label{tab:checkpoints}
\end{table}


\subsection{연구의 의문과 가설}

본 연구는 포괄적인 화면 설명이 간결한 설명보다 풍부한 웹 표현을
형성하는지에서 출발한다. Long PT가 화면의 세부 텍스트와 상태,
상호작용 요소를 더 많이 보존하므로, 첫 번째 가설은 그 이점이 PT
직후의 웹페이지 이해에서 가장 뚜렷하게 나타난다는 것이다. 다만
많은 정보를 나열하는 목표가 핵심 요소의 위치를 찾는 능력에도
유리한지는 자명하지 않으므로 이해와 그라운딩을 분리해 측정한다.

두 번째 의문은 PT에서 생긴 차이가 공통 IT 이후에도 남는가이다.
과제특화 질의응답과 명령어 supervision이 초기 표현의 차이를
덮어쓴다면 PT 직후의 우위가 최종 모델에는 이어지지 않을 수 있다.
마지막으로 동일한 행동 궤적까지 학습한 뒤 C6--C8을 비교하여,
서술 상세도의 영향이 실제 trajectory-following agent와 비웹
텍스트 밀집 이미지로의 전이까지 이어지는지 확인한다.

\subsection{평가 벤치마크}

C0--C5는 VisualWebBench의 이해와 그라운딩을 분리하여 평가하고,
WebSRC, WebClick, ScreenSpot-v2, 한국어 웹 벤치마크를 함께
사용한다. DocVQA의 ANLS와 ChartQA의 relaxed accuracy로 비웹 텍스트
밀집 시각 이해 전이를 측정한다. C6--C8의 주요 에이전트 평가는
Multimodal-Mind2Web이며 cross-task, cross-website, cross-domain
분할을 각각 보고한다. 좌표 출력은 모든 평가에서 0--999 규약으로
해석한다.

\subsection{주요 결과 (진행 중)}

\begin{table*}[t]
\centering
\small
\begin{tabular}{lccccccc}
\toprule
& \multicolumn{3}{c}{Grounding} & \multicolumn{3}{c}{Web understanding} & \multicolumn{1}{c}{Gen.} \\
\cmidrule(lr){2-4}\cmidrule(lr){5-7}\cmidrule(lr){8-8}
Pretraining supervision & Screenspot$^{\dagger}$ & GroundUI$^{\dagger}$ & VisualWebBench-G & VWB-U & WebSRC & Korean Web & DocVQA \\
\midrule
Grounding-only & 70.83 & 69.90 & 50.19 & -- & -- & -- & -- \\
Info-L  & 81.92 & 76.00 & 68.99 & -- & -- & -- & -- \\
Info-M & 82.86 & 80.30 & 62.21 & -- & -- & -- & -- \\
Info-H & \textbf{83.25} & \textbf{80.60} & \textbf{75.78} & -- & -- & -- & -- \\
\bottomrule
\end{tabular}
\caption{\textbf{Description coverage에 따른 pretraining 결과.}
Info-L/M/H는 각각 $\bar\rho=0.13/0.44/0.98$의 description coverage 조건이다.
모든 조건은 동일한 스크린샷 집합과 grounding 데이터를 공유하며,
description coverage만 다르다. SS-v2는 ScreenSpot-v2,
VWB-G/U는 VisualWebBench의 grounding/understanding을 가리킨다.
$^{\dagger}$ mobile과 desktop split의 평균.}
\label{tab:main}
\end{table*}
\begin{table}[t]
\centering
\scriptsize
\setlength{\tabcolsep}{3pt}
\begin{tabularx}{\columnwidth}{@{}Xccc@{}}
\toprule
모델 & Task & Website & Domain \\
\midrule
Grounding-only & -- & -- & -- \\
Concise & -- & -- & -- \\
Balanced & -- & -- & -- \\
Comprehensive & -- & -- & -- \\
\bottomrule
\end{tabularx}
\caption{\textbf{동일한 post-training을 마친 네 조건의
Multimodal-Mind2Web 결과.} 열은 각각 cross-task, cross-website,
cross-domain 분할을 나타낸다.}
\label{tab:agent}
\end{table}


하나의 종합 점수만으로 결론을 선택하지 않는다. 예를 들어 Long
caption은 이해 성능에는 유리하지만 정밀 그라운딩에는 불리할 수
있고, PT 직후의 차이가 IT 이후 사라질 수도 있다. 따라서 C1--C2,
C4--C5, C7--C8을 각각 비교하고 C0, C3, C6으로 PT 자체의 한계
효과를 측정한다. 벤치마크가 허용하는 경우 paired bootstrap
confidence interval을 보고한다. 다만 학습 seed가 하나이므로 통계적
주장은 run 간 분산이 아니라 평가 표본의 불확실성으로 제한한다.
